\scanpage{393}\mspage{307}

(qui existe dans $M$, puisque $M$ satisfait à $L_{\pi_0}$). Alors \deleted{\unsure{les}} $(x_{J'})_{J'\in I}$ forme un système inductif dans $M$, et $I\longrightarrow M$ commute aux limites inductives grdes devant $\pi_0$ et de cardinal $\leq\pi$.

Si $M$ est $\pi_0$-accessible, et les $x_j$ ($j\in\mathrm{Ob}\,\unsure{J}$) sont dans $M_\pi$, alors $(x_{J'})_{J'\in I}$ est un système inductif $\pi_0$-$\pi$-adapté dans $M$.

Tout ceci est \unsure{ess$^{\mathrm{t}}$.} trivial, et donne

\uline{Corollaire} 4.13.14 (de 4.13.1\unsure{2|7|3}) : sous les conditions de 4.13.12, si $\pi^{\pi_0}=\pi$, on peut pour \unsure{tt} $x\in\mathrm{Ob}\,M$ trouver une représentation $x\simeq\underset{j\in I}{\underset{\longrightarrow}{\lim}}\,x_j$, où $(x_i)_{i\in I}$ est un système inductif $\pi_0$-$\pi$-adapté avec $I$ \uline{ens. ordonné}.

\deleted{\unsure{Et}} De plus, si $\pi'>\pi$ et si $x\in M_{\pi'}$, on peut prendre $I$ tel qu'on ait

\[
\mathrm{card}\,I\leq{\pi'}^{\pi}.
\]

Ce dernier complément est immédiat,

% archive sheet number: 392
